\PassOptionsToPackage{svgnames}{xcolor}
\documentclass[12pt, a4paper]{book}
\usepackage{lmodern}
\usepackage{tikz}
\usepackage{enumitem}
\usepackage{fancyhdr}
\usepackage{hyperref}
\usepackage{titlesec}
\usepackage{setspace}
\usepackage{etoolbox}
\usepackage[object=vectorian]{pgfornament}

% 裝飾團案
% https://tex.stackexchange.com/questions/32711/totally-sweet-horizontal-rules-in-latex
\newcommand{\sectionline}{%
  \noindent
  \begin{center}
  {\color{DarkViolet}
    \resizebox{0.5\linewidth}{1ex}
    {{%
    {\begin{tikzpicture}
    \node  (C) at (0,0) {};
    \node (D) at (9,0) {};
    \path (C) to [ornament=85] (D);
    \end{tikzpicture}}}}}%
    \end{center}
  }

\newcommand{\sectionlinetwo}[2]{%
  \nointerlineskip \vspace{.5\baselineskip}\hspace{\fill}
  {\color{#1}
    \resizebox{0.5\linewidth}{2ex}
    {{%
    {\begin{tikzpicture}
    \node  (C) at (0,0) {};
    \node (D) at (9,0) {};
    \path (C) to [ornament=#2] (D);
    \end{tikzpicture}}}}}%
    \hspace{\fill}
    \par\nointerlineskip \vspace{.5\baselineskip}
  }

% 設置引用格式
\AtBeginEnvironment{quote}{\singlespacing\small}

\setlength{\parskip}{5pt}

% 用罗马数字做节的数字
\renewcommand{\thechapter}{\Roman{chapter}.}

\renewcommand*\thesection{\arabic{section}.}
\renewcommand*\thesubsection{}

% 設置頁眉頁腳
\pagestyle{fancy}
\fancyhead{} % clear all fields
\fancyfoot{} % clear all fields

\fancyhead[LE, RO]{\thepage}
\renewcommand\chaptermark[1]{\markboth{#1}{}}
\renewcommand\sectionmark[1]{\markright{ \ #1}}
\renewcommand{\headrulewidth}{0pt}
\fancyhead[RO, RE]{\textit{\rightmark}}
\fancyhead[LE, LO]{\textit{\leftmark}}
\fancyfoot[LE, RO]{\thepage}

% 設置章節標題格式
\titleformat{\chapter}[block]{\centering\Large\bfseries}{\thechapter}{1pt}{\Large}
\titleformat{\section}[block]{\centering\large\bfseries}{\thesection}{1pt}{\large}
\titleformat{\subsection}[block]{\centering\large\bfseries}{\thesubsection}{1pt}{\large}


\title{\textbf{Trades' Unions and Strikes:\\Their Philosophy and Intention}}

\author{T. J. Dunning}
\date{1860}

\begin{document}

     \maketitle

    \begin{center}
        {\LARGE TRADES' UNIONS AND STRIKES}

        \vspace{2em}

        \sectionlinetwo{black}{88}
    \end{center}

    \vspace{2em}

    {\noindent \textsc{The} object of this pamphlet is to obtain a
    hearing on the subject of ``Trades' Unions'' and ``Strikes'' from
    that class who, from their position in society, can only obtain a
    knowledge of it from the evidence of others. In setting forth the
    principle which calls Trade Societies into existence, it is not
    intended to defend all that has been done by Trades' Unions, any
    more than an advocate of religion would attempt to defend all that
    has been done in its name by religious professors. Trades'
    Unionists, generally, are neither writers, nor, to any great extent,
    able to employ writers in their behalf, nor, indeed, to the limited
    extent that they are able, do they at present see the necessity of
    doing so; while, on the contrary, those opposed to them., from their
    wealth and position in society, form a great part of the very class
    by whom writers are directly and indirectly supported;---hence,
    while the press---newspaper and periodical---teems with
    denunciations against them, scarcely a voice, directly from
    themselves, is heard through the press in their behalf Indeed, it
    has been felt by impartial men to be the one great difficulty in the
    way of understanding the subject, that all they can get hold of on
    the matter---except the merely local statements relative to a
    particular strike or dispute, which may then be pending---is
    exclusively on one side. As far as our limited means and power will
    enable us, it is intended by this pamphlet to supply that
    deficiency, by a statement direct from a body of working men, by
    their accredited agent, of their views on the subject.}

    The principal argument against Trades' Unions is professed to be
    derived from the alleged ignorance of the working classes of
    Political Economy. As this term is used by many without having a
    precise idea of what it is, and by different persons in a different
    sense, we will commence by very briefly stating the meaning which we
    attach to it: Political Economy is the science which treats of the
    principle of exchange in its operation in the production of wealth,
    and its operation and influence on civilised society.

    Man is said to be an animal---in contradistinction from all others
   ---that makes exchanges. And it is by the principle of exchange
    that all the transactions of civilised life are carried on, by
    which, indeed, civilised society exists and is held together. But if
    man be an animal that makes exchanges, he undoubtedly retains, in
    common with the carnivora to which he in part belongs, instincts of
    a predatory character, and might be defined to be an animal that,
    besides making exchanges, also steals. The act of exchange embraces
    every conceivable form in which the subject can be treated, and,
    indeed, is that out of which the science itself is evolved. It is
    only necessary, therefore, by analysing this act, to illustrate what
    we understand by the term Political Economy. By the act of exchange,
    the four things essential to a civilised community are supposed,
    without which, indeed, it could not exist, namely.
    \textsc{Labour}---\textsc{Capital}---the right of \textsc{Individual
    Appropriation}, or \textsc{Property}---and \textsc{Security}.
    \textit{Labour}, to produce the thing exchanged; \textit{Capital},
    in the accumulated labour, which results in the article exchanged;
    the recognition of \textit{the right of Individual Appropriation},
    or property, by which the act of exchange takes place, as
    distinguished from the act of taking or keeping possession of the
    article by force; and \textit{Security}, which ensures the peaceful
    continuation of a succession of exchanges to all who have anything
    to dispose of: the result of the whole, in civilised life, being the
    four great classes into which society is divided. Labour yielding
    \textit{wages} giving support to all who labour (using the term in
    its most extensive sense); capital giving \textit{profit}---the
    revenue of those who employ capital. The right of individual
    appropriation, or property, as applied to the land, which, without
    remark, we take as we find it, gives \textit{rent}, which gives
    revenue to landlords. And security supposes \textit{taxes}, which
    support those who defend the country, administer its laws, and
    tax-receivers in general. This last strongly illustrates what is
    very little noticed, if at all, by political economists---namely,
    the predatory instincts of mankind, the omission of which often
    leads to very great error in the application of rules of political
    economy, and to none more than to their application to Trades'
    Unions. And yet these predatory instincts are most unmistakably
    indicated and provided against by this element of exchange,
    security, in the absence of which it certainly could not go on, as
    in communities it requires laws with strong force to support those
    who administer them and enforce their execution, without which the
    community could not exist; and in separate nations, armies and
    fleets for defence, and, as a stronger illustration, for attack. On
    these four elements of the principle of exchange, without either of
    which it could not exist, it will be seen the whole system of
    civilised society rests. We have thus particularised each element,
    and its respective result, to show that those who live by labour
    are, by \textit{right}, entitled to their revenue by the simple
    action of the principle of exchange equally with the other three
    great divisions of society, as opposed to the confused idea of wages
    being derived from the relation of master and servant, involving a
    sort of benevolent patronage on the one side, and obligation on the
    other. And it will be our object to show, that it is only by trade
    combinations that the \textit{free} exercise of this principle can
    be secured to working men.

    Our business, however, is only with the action of the principle of
    exchange upon wages. It is usual, with a certain class of political
    economists, of whom A. Black, Esq., M.P., may be considered a fair
    exponent, to lecture the working classes on the principles of this
    science. Accordingly, in his recent lecture in Edinburgh, he devotes
    five, out of the twenty-two pages it comprises, to the consideration
    of the value of commodities---labour being one in its nature like
    the rest---which is determined by demand and supply; the pernicious
    result of monopoly; of capital and labour, and of each being
    necessary for the employment of the other; that capital will not be
    employed unless a profit can be obtained by its employment; that
    when the profit of capital is high in any trade or manufacture, so
    many will be tempted to embark in it that this high profit will soon
    be reduced to its ordinary level; and the same with high wages,
    except they proceed from great skill, which, from the nature of
    things, must always be rare; and, lastly, if a capitalist has
    embarked his capital upon the supposition that he has to pay a
    certain sum for wages, and he is obliged to pay more, his profit
    will be reduced by that amount, which will be very likely to cause
    him to withdraw from the business in which he has engaged, and, of
    course, cease to pay wages, which should be a lesson to workmen not
    to demand too high wages.

    With due respect for Mr. Black, who doubtless was actuated by a
    sincere desire to promote what he believed to be the true interests
    of the working classes, we think, if, instead of puzzling himself
    and his readers by elaborately explaining the above propositions,
    which nobody is likely to dispute, and which are, in their
    enunciation, like that of two and two making four, already at their
    plainest terms, he had endeavoured to show how it was that these
    things, which almost everybody knows, failed to convince the working
    classes that Trades' Unions are unnecessary, and Strikes always
    pernicious, and shown how they were in error: he would then have
    conferred a real benefit on the working classes. It is a great
    mistake, in gentlemen of Mr. Black's position, to suppose that
    working men are practically ignorant of the truths of political
    economy. With it as a science they may be ignorant; indeed, some of
    them may not even know that there is such a science. But these
    gentlemen seem to forget that political economy is only the
    compilation into a code of rules the practice of all men in all ages
    who have bought and sold; have been engaged in commercial
    transactions; or in any way have made use of the principle of
    exchange. Every costermonger is obliged to regulate what he does in
    buying and selling by the law of demand and supply, which if
    explained to him as being a part of the science of political
    economy, he would be as much surprised that he had, without knowing
    it, been practising the science all his life, as the honest
    gentleman in the French comedy was, on his being taught composition,
    that he had been talking prose all his life without knowing it.
    Except in the innocent dulness of the latter, which would certainly
    not belong to the costermonger, the cases are exactly parallel.

    We therefore will endeavour to show why the above propositions fail
    in convincing working men that they are doing wrong in entering into
    combination. Not that these propositions arc doubted, but because
    implicit belief in them is perfectly consistent with the propriety
    of Trade Combinations. We shall therefore proceed to consider ---

    \begin{enumerate}[label=\Roman*., leftmargin=2.9cm]
        \item---Wages, and what determines their value.
        \item---Trade Societies, for the protection of wages.
        \item---The means used by them for that purpose.
    \end{enumerate}

    \chapter{--- WAGES.} \label{WAGES}

It is superfluous to say that the price of labour, like that of
everything else, is determined by the quantity or supply of it
\textit{permanently} in the market; when the supply of it
\textit{permanently} much exceeds ils demand, nothing can prevent the
reduction of wages; and, conversely, when the demand for it permanently
much exceeds its supply, nothing can prevent their rise. In these two
extreme points all contention is hopeless. No Trade Society on the one
hand, however well organised, can, or ever did, prevent the fall of
wages in the first case; nor, in the other, can, or ever did, the most
stringent legislative enactment, of which there have been many
instances, prevent their rise. Trade Societies, however, rarely meddle
with these two extremes. Leaving them, we come to the intermediate
states that admit the operation of Trade Societies; and, indeed, which
call them into existence.

In all exchanges, besides the adjustment which takes place in them by
the operation of demand and supply, there is always, from the predatory
instinct inherent in the very nature of man, noticed above, a desire, on
one side or the other, to take advantage of the necessities either of
the buyer or seller; and, in proportion as these necessities are
immediate and pressing, to press that advantage accordingly. Take, for
example, an estate for sale, of which it is known that, come what will,
it must be sold in a fortnight for ready money, or what is deemed as
good as ready money. If the estate be large, the seller, under these
circumstances of necessity, will be sure to lose some thousands of
pounds. In wages, besides the rate of wages, which results from the
demand for it in proportion to its supply, there is a lower rate which
may be the result of the necessities of the workman. For example, in
those trades where there is what is called the ``Sweating System''
practised, the fair result of the demand and supply rate of wages is
represented by the amount received by the ``Sweater;''\footnote{A
``Sweater'' is one who takes out work to do, at the usual rate of wages,
and who gets it done by others at a lower price; the difference, which
is his profit, being ``sweated'' out of those who execute the work.
Hence the term ``Sweater.''} perhaps, however, his demand, to enable him
to get the work, may be somewhat less. He, however, gets the work done
lower still; which is his profit, by employing those whose necessities
compel them to work for his price, or starve---some of these, perhaps,
may be first-rate workmen, whose habits of intemperance have reduced
them to the lowest state of destitution, and who are glad to work for
any price. There is, therefore, a wide difference between the demand and
supply rate of wages, or the rate which the fair operation of exchange
would give, if the buyer and seller of it were on equal terms, and that
which is, or would be, compelled, if the employer dealt with each man
``singly,'' and obtained the reduction which his necessities might
dictate.

The supply of labour in a trade may be greater at one time than another
--- indeed for a time far exceed the demand for it; and yet that state
not be its permanent condition. In other words the trade might be for a
good while ``slack.'' If a reduction of wages took place during this
period, this reduction would be very likely to remain when the trade got
``busy,'' the supply of hands not being then greater than the demand;
and yet the temporary ``slackness'' in itself would by no means be
adequate to compel this reduction, and ought not to have produced it.
The undue competition among employers, bidding each under the other,
intending to make up the difference out of the wages of their men,
---not from any compelling necessity in the trade, but from mere
rivalry, ---would inevitably cause a reduction, which, if not checked,
might extend to the whole; not again from any compelling necessity
arising from over-supply, but from taking advantage of the immediate
necessities of their men being greater than their own. Hence the
formation of Trade Societies, which become a necessity, to adjust the
bargain for the sale and purchase of labour.

In all bargains, the buyer wishes to buy as cheap, and the seller to
sell as dear, as he can; but their interests, all being exchangers, and
each, from the highest to the lowest, depending upon this principle
---exchange--- for his position in life and even for his daily
sustenance; all things being equal, their position is not one of
opposition, but of mutual interest; and neither the one in wishing to
get as much, nor the other in wishing to pay as little as possible, can
injure the other. But if, as alluded to above, either party possess an
advantage over the other in the bargain, this position of mutual
interest is changed into one of opposition, in which the weaker party is
sure to be deprived of some proportion of what is justly due to him. In
this position as bargainers for the sale and purchase of labour, stand
the employer and employed. Singly the employer can stand out longer in
the bargain than the journeyman; and as he who can stand out longest in
the bargain will be sure to command his own terms, the workmen combine
to put themselves on something like an equality in the bargain for the
sale of their labour with their employers. This is the
\textit{rationale} of Trade Societies, which is very clearly indicated
by Adam Smith in his ``Wealth of Nations.''
    \chapter{---TRADE SOCIETIES.} \label{TRADE_SOCIETIES}

The object intended is carried out by providing a fund for the support
of its members when out of employ, for a certain number of weeks in the
year. This is the usual and regular way in which the labour of the
members of a Trade Society is protected, that the man's present
necessities may not compel him to take less than the wages which the
demand and supply of labour in the trade have previously adjured;
strikes, which we shall hereafter notice, being only resorted to on
extraordinary, and, generally speaking, most unusual occasions.

As the object of combinations is thus to correct the disadvantage of
position as bargainers with their employers, and as the very element of
their existence is public spirit, being held together only by a
principle of honour to support each other for a perfectly legitimate
purpose, they always comprise not only the best workmen, but best men in
a moral sense that are to be found in the trade. This is stated, simply
because it is a fact, which has been acknowledged by those who are
inimical to trade combinations. As the social position of workmen
depends entirely upon the wages they obtain, it is felt by the majority
of them to be a sacred duty to adopt this means of protecting their
wages. There is no state so abject, and no tyranny so hard to be
endured, as that which results from the seller being completely at the
mercy of the buyer. It is to avoid this state so detrimental to
themselves, and so advantageous to the employers ---at least so thought
by many of them, ---that working men combine, and that their
combinations, according to the ``Edinburgh'' reviewer, contain ``about
600,000 members, commanding a fund of £300,000.'' How this fund is
disposed of we shall hereafter allude to. This is the reason, at which
neither he nor any man need ``for the twentieth time wonder,'' that
induces ``intelligent workmen'' to form themselves into Trades' Unions;
and not, as stated by the same reviewer, after about four pages of vague
denunciation and foolish wonder, to ``intimidate,'' ``to stint the
action of superior physical strength, moral industry, or intelligent
skill,'' or ``to apply the fallacies of the protective system to
labour;'' their object, on the contrary, being to ensure the freedom of
the principle of EXCHANGE with regard to labour, by, as before observed,
putting the workman in something like an equal position in the bargain
for sale of his labour with his employer. If in doing this they
sometimes---err as who does not?---such error does not invalidate the
propriety of such a course.

The same reviewer, in a page which, though placed at the commencement,
seems to have been written as a kind of summing-up of the whole, asks
his readers: ``Can it be possible, then, that in Great Britain, too,
there is an \textit{imperium} in \textit{imperio} as oppressive to those
whom it comprehends as any secret society on the Continent, enforcing
mysterious laws and arbitrary obligations by the hands of irresponsible
authorities, in defiance of the great natural laws of human society?
Recent events have prepared our readers to take some interest in the
inquiry whether the fact is so or not. When they have learned as much as
can be known of the Trades' Unions of England, they may judge for
themselves whether any existing government in Europe from Constantinople
to St. Petersburg or Paris would venture to exercise so stringent a rule
over its subjects as a large proportion of our working men submit to
from men of their own order. Under the system of irresponsible
government, which strives to subject the labour market to its dominion,
the characteristic freedom of the English citizen is lost. Hundreds of
thousands of our countrymen are unable to say that each man's house is
his castle; that his labour and its rewards are his own property, and
that `live and let live' is the rule of the society in which he
dwells.'' ---P. 528.

The sum and substance of the above, divested of its gloss, is that
Trades' Unions are, 1st, secret; 2nd, irresponsible; and 3rd, coercive.

\section{---SECRET.}

With respect to the ``mysterious laws and arbitrary obligations,'' we
have only to say that the proceedings of Trades' Unions are no more
secret than those of any other societies, to which working men and
others belong. When the reviewer talks of ``Continental nations from the
Mediterranean to the Volga, and from the Black Sea to the Channel, being
honeycombed with secret societies,'' and insinuates that it is a mistake
in Englishmen to congratulate themselves that everything with them is
``open and above board,'' and that they have no ``spies,'' ``enforced
dumbness,'' and ``\textit{surveillance},'' while Trades' Unions exist
among them; he insinuates what is false. The insinuation, for he does
not venture to state it openly, without doubt, gives strength to the
alleged tyranny, but it is untrue. That the proceedings of Trades'
Unions have been unknown is true, but they have only been unknown
because they have been considered by the general public, and therefore
by the press, too obscure and insignificant to deserve attention.
Nothing would serve the cause of the working classes more than a
continual ventilation of their proceedings by the public press; it
would tend to procure for them a greater meed of justice than they have
hitherto found, and, what is better, it would tend greatly to preserve
them from error. There is not a more salutary check to injury from your
own proceedings, as well as from the proceedings of your opponents, than
publicity through the press. This, however, has undergone a change for
the better, during the last few years. In the ``builders' dispute,'' as
far as the men's proceedings are concerned, it has ceased altogether;
the secresy being entirely on the side of the Masters' Association.

We have now before us 227 quarto pages, double columns, for the most
part in the \textit{small} newspaper type, of the reports of the
proceedings of the men, leading articles, \&c., on this dispute, cut
from the newspapers, chiefly from the \textit{Times, while that paper
reported them}, and when it ceased to do so, from the \textit{Morning
Advertiser}, and other papers. This, in the usual type (not the largest),
would make a quarto volume of 671 pages. Where, then, is the secresy?
The masters' proceedings, in their Association, have been secret, but
not the men's, in their ``Trades' Union.''

There are no ``secret societies'' as ``Trades' Unions.'' The attempt to
make them so, in 1834, was utterly abortive, from the refusal of the
trades generally to become so. When it -was illegal to combine, they
were necessarily secret, not from any inclination, but from the force of
the adverse law, which compelled secresy. This enforced secresy
generated its usual train of evils, and in some cases the administration
of oaths to members. When the laws against combinations were repealed,
these practices, and the secresy out of which they grew, ceased, which
ought to have been known to the reviewer, before he attempted to write
on the subject.

The ordinary proceedings of a Trade Society are of course too
uninteresting for public attention. But in disputes, which in the case
of the builders, have occupied so much of the attention of the public,
whose fault has it ever been that they have not obtained the same
publicity? Certainly not the men's. During the strike of the masons at
the building of the two Houses of Parliament in 1841, the greatest
difficulty was experienced by the men in getting their statements into
the newspapers; in fact, they were obliged to resort to placards and
printed bills, because the newspaper press refused to publish them. The
masters found immediate publication in the first-class newspapers for
their statements, but the men could not. The present writer had the
greatest difficulty in getting a report of an investigation into a case
alleged by the men and denied by the masters, relative to this dispute,
into a popular newspaper, that now would publish such a report
immediately. He with his colleague had to wait upon the editor, who had
previously refused insertion; and at last it was inserted---as he
believes---through the intervention of one who was influential in the
management of the paper. We go into this detail to show that it has
never been the fault of Trades' Unions that their proceedings have been
unpublished to the world. But it is said the Committees are secret. It
is true their proceedings are not published, but does that constitute a
secret Committee? If so, then all meetings of men appointed to transact
business, from the Cabinet Council downwards, including all railway and
parochial Committees, are secret conclaves. But these are not secret,
because the names are known of those composing them; the first-named to
all the nation, and the rest to all whom they concern. And, in like
manner, the names of those who form Trade Committees are equally known
to those who appoint them; and, therefore, are not secret.

\section{---IRRESPONSIBLE.}

The ``irresponsible authorities,'' by whom we suppose are meant the
President, Committee-men, and Secretary, of each and every Trades'
Union---if the reviewer was so ignorant of such matters as not to know
it before---these are invariably elected by the Trades' Union to which
they belong, to which they are responsible for what they do during their
term of office, and who are not slow to find fault and censure them if
they do anything displeasing to the body they represent. If he did not
know this, the commonest sense of propriety and regard for truth ought
to have restrained him from writing on the subject until he had obtained
this information. But then there could have been no ``slashing article''
the dull prose of truth would have prevented it. As the responsibility
of the ``authorities'' of Trade Societies will be further shown as we
proceed, there is no need here of saying more.

As to the ``characteristic freedom of the English artizan being lost''
by this ``system of irresponsible government,'' the reviewer
unconsciously, from his ignorance of the subject, has, by the
\textit{reductio ad absurdum}, refuted himself. He describes a state of
things which has no existence. We should no more have believed this
passage about working men, not knowing whether their ``labour and its
rewards were their own property,'' and the rest of it, applied to
ourselves as workmen, had we not been told, than we should have known
that the French writer meant Shakspere when he called him the ``divine
Williams.'' We have shown this passage to several of the ``hundreds and
thousands of our countrymen'' here described, who have one and all, with
much laughter, in the most expressive vernacular, and with vigorous
expletives, denounced the reviewer as the ``biggest of the unwise,'' and
the greatest of the perverters of truth, that they ever heard of. When
any matter, like the present, has reached its last point of absurdity,
there is no more to be said. We therefore take leave of it, especially
as our object is not to criticise the ignorance of the reviewer, but to
state the truth; and, having disposed of the secresy and
irresponsibility---two such powerful adjuncts to the alleged coercion
``which strives to subject the labour market to its dominion''---we
will now proceed to consider the alleged tyranny or coercion itself.

\section{---COERCIVE.}

As the ``coercion,'' apart from acts of intimidation and violence, said
to be exercised by Trades' Unions, has been the theme of accusation,
invective, censure, or expostulation, according to the mood of the
speaker or writer, with the unanimity of a chorus, from Lord Brougham
down to Mr. Buxton, we hope to be excused if we enter somewhat largely
into the subject. It is not only alleged to indnce men to belong to the
Union, but to exist everywhere, in all that Trades' Societies do, or
intend to do. Now, it is the tyranny of the majority over the minority;
then, it is the tyranny of the leaders; or, again, a general undefined
tyranny which pervades the whole body; but always tyranny and coercion.
According to the ``Edinburgh Review'' the tyranny is general, undefined,
but pervading everywhere. The \textit{Times}, Dec. 9, 1859, gives it to
the leaders in the following terms:---

\begin{quote}
    % 简单粗暴第在引用段落中实现缩进……
    \hspace{2em} They strike at the bidding of men who deem it necessary
    to assert a control over the destinies of trade by making a terrible
    example of employers and employed as often as their secret and
    imperious commands are disobeyed. Nor does their fault consist in
    taking too much, but rather in taking too little, interest in their
    own affairs---in permitting themselves to be made mere tools and
    instruments of the will of others in matters in which it is, above
    all things, incumbent on them to judge for themselves, and in
    surrendering their natural liberty, tamely and ingloriously, into
    the hands of an ignorant and inquisitorial despotism.
\end{quote}

Mr. Black, taking up the subject from the ``Edinburgh Review,'' the said
reviewer having used the word ``excommunication,'' says:---

\begin{quote}
    \hspace{2em} But it is not by illegal violence alone that personal
    liberty may be encroached upon; there are other means of
    intimidation that may be as effectual, or even more so. The Pope, by
    the threat of excommunication, was able to control the most powerful
    sovereigns. A similar threat is one of the most effective
    instruments in the hands of the committees of the Unions. When they
    order a strike, he would be a bold man who would preach the right of
    private judgment, and act upon it in opposition to the commands of
    the secret committee; he would not only be shunned by all his
    companions, but reproached and denounced as a snob, a nobstick, as
    unfair, and insulted in other opprobrious terms. By cajolery,
    molestation, and threatening, the combined leaders strike a general
    terror into the whole community.---P. 21.
\end{quote}

The ``Quarterly'' reviewer, having more sense, does not harp so much on
this string.

To Mr. Black is due the merit of being definite; he not only states what
is done, but the manner of execution. And we have no doubt but that he
has, in the case he has supposed, stated the true impression on his
mind, as derived from the ``Edinburgh Review,'' and the statement of
various speakers and writers on the same side, who are considered
authorities on the subject. We will, therefore, analyse it. The
``excommunication'' is relative to a supposed strike, and he assumes it
to be complete and effectual. Now, how can a man be ejected or driven
from a community, with ``reproaches and denunciations,'' ``shunned even
by his friends,'' unless by the active co-operation of the community
itself? and how could such co-operation exist by the community, unless
such `` excommunication'' had its full and hearty approval? No cajolery
can make men, and especially working men, do what they have no mind to,
and, if the community are willing to eject one of their body so
vigorously as here described, with ``insult'' and ``opprobrious'' terms,
there needs no ``cajolery'' to induce them to do so. As to the
committee, who are as one to fifty, or more, attempting to intimidate
the general body to do such a thing by threats, it is too preposterous
for argument. Mr. Black here describes the tyranny of the majority over
the minority; and not, as he states, the tyranny of its leaders over the
general body. His own case proves that he is entirely mistaken in this
respect. It may be, in such a case, that the leaders are compelled to
act in obedience to the general voice, as has been the fact, hundreds of
times, with more important bodies than Trade Societies. So far from the
fact being, as the Times and Mr. Black suppose, that the Committees of
Trades' Unions are the agitators and leaders in strikes; it oftener
happens than not that they prevent them, by soothing down the angry
feelings excited by what the members of their Societies have considered
unjust treatment, from foremen or employers. There, however, exists some
cause for this error. When men have been asked, by employers, why they
have struck or have refused to work for a certain price, they too often,
from the fear of compromising themselves, lay it all to ``their mates,''
or to ``the Society,'' or to its ``Committee,'' when perhaps some of
those very men have been foremost in setting forth to the meeting of
their general body, which unanimously agreed the strike should take
place, the absolute necessity of so doing. But, although this cause for
error does sometimes exist, it is not the less a great mistake to
suppose that the working class are the mere ``tools and instruments of
the will of others,'' or that they ``surrender their natural liberty''
into the hands of their leaders. Those leaders, as before observed, are
in all cases elected by the general body, to carry out its views, and
therefore cannot be secret. Nor can they do anything against the will of
that body. If they were to attempt to do so, they would soon cease to be
its leaders, and, consequently, could not exercise ``despotic'' rule
over it.

But it may be asserted that, granting this, the rule of Trades' Unions
is nevertheless the tyranny of the majority over the minority. To which
we reply, where is it that this kind of tyranny does not exist? Mr.
Black tells us that there are no ``Trades' Unions'' or ``Strikes'' in
the United States of America, and we believe he is nearly correct.
Where, according to M. de Tocqueville, is the tyranny of the majority
over the minority greater than in those states? To come nearer home,
what is the ``world's dread laugh'' but the tyranny of the majority?
What is ``public opinion,'' so potent, and so fatal to those who stand
in opposition to it, but the voice of the majority against the minority?
who never fail to denounce it as tyranny; and who generally get laughed
at for their folly by the same ``public.'' If this be a crime in Trades'
Unions, it is a crime of which the \textit{élite} of English society are
equally guilty; and not only the \textit{élite} but every grade
downwards from it to the working classes. Out of many instances of the
way in which it sometimes acts that might be given, we only need mention
the following:---Mr. Burke, after his expression of opinion on the
French Revolution, was, to use his own words, ``excommunicated by his
party,'' lost his seat in Parliament, and for his whole life after was
reduced to political insignificance. Lord Chancellor Yorke, who took
office at the earnest request of the King (George III.) from being
supposed by such taking of office to have betrayed his party, was
proscribed by that party; which he took so much to heart, that he lost
his reason, and died by his own hand.

Sir Robert Peel was considered, by the party to which he had for years
belonged, to have betrayed it, though it was all the while known that he
had a perfect right to change his opinion; and the event, during his
life, proved that it was most beneficial to the nation that he had done
so---yet with his party he lost caste; and though he had the approbation
of the general public, his public life was embittered to the end---even
when the aggregate majority was with him---by the most bitter attacks
that ever assailed a public man; against which the approbation of the
general majority was of no avail, such approbation being expressly
because he had shattered to pieces the favourite policy of his former
party; or, to use its owm words, betrayed it.

Messrs. Bright and Cobden, during the late war with Russia, thought
proper on the peace question to express what was known previously to be
their sincere convictions, and which they had an undoubted right to
express, but which were opposed to what was at that time considered to
be the public weal; in short, to the almost wishing success to the enemy
with whom the country was at war. All their past services were
forgotten, they both lost their seats in Parliament, and, during that
war, their influence as public men; and would again under the same
circumstances suffer the same loss, supposing they should, which is very
doubtful, ever regain the confidence which the public formerly reposed
in them. Pausing for a moment on this phase of the subject: do either of
these reviewers, or any one else, require to be told, that a state of
combination, whether it be of a political party or a Trades' Union,
supposes a united action which is inconsistent with dissent?---or that
there are men everywhere, who, after the whole of the fellows---
themselves included---have agreed to a certain movement, have frustrated
it by an opposite course of action? Have not many such been denounced in
both reviews as having ``betrayed their party,'' and ``their cause?''
Without doubt such men have a right to do as they like, and to be
protected from violence and insult; but it is too much to expect that
their fellows should receive them with affection or regard. Do what you
will, talk how you like, such men will always be considered by their
party, whether it be of the highest or lowest sections of the community,
as infamous, and treated accordingly by the majority, whose cause they
have betrayed. If this be ``tyranny,'' then it will never be purged from
the working or any other class of society, while men are human. When, if
ever, they become angelic, it may be different; but then there will be
no traitors.

But to return. What is the end and aim of the ``Quarterly and Edinburgh
Reviews'' as far as politics are concerned, but to excite the public
opinion against their opponents?---and the greatest success they hope to
achieve is to obtain the general voice against those whom they denounce,
or, in short, to create this ``coercion'' or ``tyranny'' of the majority
against them. What is the ``Edinburgh Review'' now doing, but, Tight or
wrong, endeavouring to excite public opinion against Trades' Unions?
What is tbe object of Mr. Black's lecture and Mr. Buxton's letter---both
compiled from the two reviews in question---but to do the same thing?
But we shall be told they all think they are doing right. Do they? We
should be glad to be informed who does not when thus employed? But, if
they think they are doing right in thus invoking public opinion, can
there be a more profound acknowledgment of its power, and of submission
to its voice, than such invocation? It is the acknowledgment of a
superior power by appeals to it. If, then, all in their turn acknowledge
its power by invoking it, as a rule to which all must submit, how can
it be a degrading tyranny when it exercises its force among working men?

But what do they want public opinion to do, when all in their turn
invoke it, but to bring its force to bear against their opponents?---to
bring, if successful, what those opponents would most surely call this
tyranny of the majority against them, with all its irresistible forms of
coercion? And this is deemed by all parties not only legitimate and
proper, but in the highest degree commendable. Is it, then, only when
this formidable power produces the same result among the working
classes, that, in the opinion of this reviewer and the rest, it is
tyranny and coercion? Take the question as it is put by Lord St.
Leonards, in his letter to the \textit{Times} of December 13, 1859, and
the same result will follow---namely, that it is by the opinion of the
majority that Trades' Unions exist---who are said in these unions to
coerce all the rest---that is, the minority. His lordship, as far as we
can understand him, wishes to pat an end to Trades' Unions, and of
course describes their action as most unwise and tyrannical. He says
that, ``If a law were to be passed, placing the free workmen under such
a rule, we might expect an insurrection; and yet vast masses voluntarily
place themselves under a yoke which they would resist in the shape of
authority with all their energies.'' What are the vast masses here
alluded to, but the majority f And if they elect to place themselves
under what his lordship is pleased to term a ``yoke,'' he may rest
assured that they have their own good reasons for doing so, and do not
require, however well intentioned, the commiseration of anyone. Here, as
in the case put by Mr. Black, to suppose that the ``mandates'' of the
committees, \&c., would for an instant be regarded by the general body
comprising the Trade Society, unless those ``mandates'' had previously
had their sanction and concurrence, is too ridiculous for argument.
These officers and committees, once more, do not make laws for the
general body, they are made by the general body itself in all cases; and
if it assigns a discretionary power to its officers, it is because
it---namely, the body at large---think that the best course to pursue.
It seems to be forgotten that working men are not so ``handy'' with the
pen as lawyers' clerks, and therefore the real fact---namely, that it is
better to leave to the discretion of the committee the dealing with
cases that may arise, than to make rules concerning supposed cases which
might, in practice, be wholly inapplicable---is not believed. But,
instead, alarm is taken, or pretended to be taken, at the ``
irresponsible power'' of the committees. Never was alarm more
groundless. Only let any committee or officer act in opposition to the
wish of the general body, and, as before said, they would speedily find
themselves relieved from ``the cares of office.'' But why is all this
care for the ``vast masses?'' They have never, as we have heard,
complained of the ``yoke'' under which they are said to be bowed down by
an ``irresponsible and secret tyranny.'' No complaint of theirs has
evoked all this sympathy. Loud have been their complaints of the tyranny
of the master builders' ``lock-out,'' and large have been their
contributions to resist it. But not a word against this ``yoke,'' so
feelingly alluded to by his lordship. What is the reason of this? It
cannot be because they do not complain when they see a reason:---the
masters' ``dlock-out,'' for instance? What can be the cause of this
silence? We will tell his lordship and the chorus of commiserating
reviewers, lecturers, letter-writers, and others. It is because they,
that is, the ``vast masses,'' believe that by their organisation---``yoke''
his lordship is pleased to term it---they obtain a fair exchange for their
labour, which it would be impossible to secure singly as individuals:
in other words, that they obtain higher wages in combination than they
could if they were not; and they are right. Take the engineering trade,
in illustration. In Glasgow and Newcastle the custom is, not to belong
to the Engineers' Society; in those two places the wages are low
compared with their rate in those districts where the custom is to
belong to it. And even in those two places the men who do belong to the
society obtain a higher rate of wages than the non-society men. Such
instances might be multiplied in every trade. It is, therefore, a matter
of surprise that his lordship, knowing, as he does, what men will
endure, what risks they will undergo, for increased remuneration, which,
as he knows also, they would do for no living man,---that he should think
it extraordinary that the working classes should ``voluntarily place
themselves under'' this ``yoke,'' even at his own estimate of it, to
secure the same thing for themselves.

But all this points with increased force to the fact, that it is the
majority, the ``vast masses'' who determine to combine, and who in
combination decide for themselves what course they shall pursue.

Does his lordship, in alluding to ``tens of thousands'' of men being
suddenly deprived of work---as if they had been struck by paralysis---
allude to the builders' dispute? If so, he should remember that Avas the
act of the master builders by their ``lock-out.'' To this, however, we
shall refer elsewhere. Having thus shown that there is no ``coercion''
necessary to induce men to combine, but that of the ``public opinion'' of
the class who combine; and that trade combinations are neither secret,
irresponsible, nor their rule dependent on the will of their leaders; it
remains only to notice the charges that workmen ``coerce'' all in their
trade to belong to the trade society

Whatever may have been the practice formerly, Trades' Unions are not now
generally charged with acts of physical violence---acts which we
deprecate and abhor. On this point we are in perfect accord with both
the `` Quarterly'' and ``Edinburgh'' reviewers, and those who have
compiled speeches, lectures, and letters therefrom. Wherever such acts
occur, they are a ``heavy blow and great discouragement'' to trade
societies. Leaving these, as things which happily have now passed away,
we need only say in reference to this, that there is no other coercion
than that of the public opinion of the class to which workmen belong.
Although we know of many instances to the contrary in our own trade, to
a certain extent men in a society shop \textit{are} expected to belong
to the society, for the very natural reason, that, as it is believed
wages are kept from falling below the demand and supply rate, by the
contributions of men who belong to the society, those who do not belong
are considered in the light of those who wish to enjoy an advantage
without contributing their share to sustain it; nothing being more
certain, in the minds of the workmen, that, but for the combination, the
immediate necessities of the men would be immediately taken advantage of
in some part of the trade, which would speedily produce the same result
throughout. The predatory instinct belonging to our nature would in
these circumstances ensure this result. The men, therefore, naturally
expect that every man should pay his quota, for an advantage which he
enjoys in common with the rest. And Ave should be glad to be informed
how, in these circumstances, it is possible to be otherwise? It may be
considered---and especially by those employers who find it most
advantageous for their individual profit to deal with the men singly---
very wrong for the workmen to act thus. But it is for this very
employer's reason that the men are anxious that all should be in
combination. In fact, whatever is deemed to be the right and proper
course for the welfare of all, by the majority of a community or body of
men, and adopted by it in the aggregate, will be sure to be considered
by that community as the duty of all to support and carry out. In fact,
it is its public opinion---and how potent public opinion is need not be
repeated. And if that course involves a money payment, he who refuses or
attempts to shirk that payment will always be, let him be of what class
of society he may, ``coerced'' by its public opinion. It is so
everywhere. It may be very wrong, but if it be, men of every grade of
intellect and position in society have always done so, and, without
doubt, will continue to do so to the end of time.

But this opinion, in most places, is so universal among working men,
that in point of fact there are none to ``coerce;'' whether actually in
the society or not, scarcely a man doubts the propriety of belonging to
it. And where, as in some few districts, it is otherwise, the men do or
do not belong to the society, as they please, and consequently there is
no ``coercion.''

In our own trade those out of the society are, for the most part, either
inferior workmen, employed on inferior work at reduced wages, or those
who have belonged to it, and been erased. Of these last some left
because they did not wish to pay to it, or indeed to anything else that
they could avoid; and the rest, by far the greatest number are those who
have been erased for non-payment through their unfortunate habits of
intemperance, which left them no means of paying. The first-named, from
their inferior ability as workmen, seldom come to a society shop; the
others, when they obtain work in a society shop, generally join, or
rejoin; and, as the result of nearly twenty years' experience as
Secretary, the present writer can safely affirm that in no case have
those who have so joined, or rejoined, expressed the slightest complaint
of any ``coercion;'' on the contrary, where there has been any
expression at all, it has always been that they had been out of the
society so long, or had not joined it sooner. The practical fact is, the
chorus of assertions to the contrary notwithstanding, that the belonging
to the society is never by working men, taken generally, deemed a matter
of ``coercion'' at all, but, as a thing for the protection of their
wages from undue reduction, highly advantageous for them to do.

\subsection{UNIFORM RATE OF WAGES.}

Before discussing this topic, it is necessary to state that the
\textit{men} do not contend for an ``uniform'' rate of wages. All they
contend, or ever have contended for, is a ``\textit{minimum}'' rate of
wages, leaving the employer to pay for superior skill, or working
ability, as much more as he pleases, or the man can obtain. Indeed, as
great skill or superior working ability must, from the nature of things,
be always rare, what is termed an uniform rate can only mean what is
applicable to the general run of men, and in point of fact a
\textit{minimum} rate as regarded by the men, and an ``uniform rate'' as
regarded by the masters.

To have an ``uniform rate'' of wages is said to give an unfair
protection to the unskilled workman. The ``Quarterly'' reviewer states
it to be ``a most dangerous thing for workmen to proclaim that the idle
and unskilled shall be paid as the industrious and skilled.'' This
reviewer shows much candour throughout his remarks; he errs, however, in
this matter, from receiving his impression from only one side of the
question.

It is quite true that in most trades, where the work is paid by the day,
an uniform rate of wages is paid; but the above result, which has been
so much denounced, certainly is not, and never was, the intention of
such uniform rate or the reason why it was established.

It should here be observed that where the work in any trade is paid for
by the ``piece'' at so much per job---as, for example, among the
compositors, the type, in London, is for the most part composed at so
much per 1,000 letters---there is no uniform rate \textit{received} by
the workmen. There is generally an uniform rate, or nearly so, of the
price of the various jobs; but, as workmen of different quickness and
skill will do more or less work at the same price per job, their wages
may very materially differ in amount.

But in work by the day there is generally an uniform rate, which is
adjusted in the first instance by the principle of demand and supply,
which operates like the circulation of the blood without anyone scarcely
being aware of it, and the result is a settled rate of wages, which
becomes recognised by both employer and employed as the standard of
wages for the general run of workmen in the trade. And it must always be
borne in mind that a rate of wages, in any trade, never becomes what is
termed the uniform or standard rate, unless it has had the full
consent---very often by special agreement---of the employers of that
trade. It is quite true that some men are able to earn, and do earn,
more than others. As the uniform rate is generally made for the general
run of men, some men will be worth more and some worth less than its
amount. Then the question returns, why not pay the one more and the
other less? This is sometimes done, when, of course, the rate ceases to
be -uniform. But where the uniformity, as a general practice, is kept
up, the reason usually is---that the employer likes to reap the benefit
arising from the man being worth more, and is also afraid that if he
increased this man's wages, he would probably be called upon to raise
the others; and the men, en the other hand, would be against the man who
was not worth the regular rate taking less, lest the others might be
reduced to it. And it must be admitted that, so prone are men to take
every advantage in their own favour, the apprehension of each class
respecting the other is very likely to be correct. But it will be seen
in the sequel that this uniform rate operates in a way which is anything
but a protection to the idle and unskilled. Employers never fail to
retain in their employment the best workmen, and to discharge, at the
first opportunity, those whom they consider inferior. In the course of
time, by this process, they will have in their employ few but what may
be termed superior workmen, to whom they will pay by the operation of
this ``uniform rate'' of wages only the wages of the general run of
workmen. Digressing for a moment, as far as we can learn, this was
actually the case with the master builders, previous to their
``lock-out.'' They had in their employ the \textit{élite} as workmen of
their respective trades. Hence the great loss they have sustained by
their `` lock-out,'' which may be thus estimated. The number locked out
was under 11,000,\footnote{This is the number stated by Mr. Potter and
the Building Trades Conference. Wishing to be correct, Mr. Wales,
Secretary of the Master Builders' Association, No. 8, Great St. Helens,
was written to, requesting him, as a great favour, to state his estimate
of the total number of men locked out, skilled and unskilled. As Mr.
Wales did not reply, we have no means of collating the above number with
what that of the employers might be. However, from various
circumstances, we have good reason to believe the above number to be
correct, and our estimate of the loss very much understated.} skilled
and unskilled. The master builders, when there was reason to believe
they still wanted more hands, stated, in a return by their secretary to
the \textit{Weekly Mail}, of Nov. 13, 1859, the number in their employ
to be upwards of 16,000 (up to Nov. 5, 1859), to whom they would
certainly not be paying less per man than they were previously paying
the 11,000 for doing the same quantity of work, involving a loss, if wx
suppose the average wages of each to be 25s. per week, of upwards of
6,000 a week in wages alone. But to return. Inferior workmen, it will be
seen, by the operation of this ``uniform rate'' are not protected, its
result being to give the employer the benefit, as described, of having
superior workmen at the ordinary rate of wages. Indeed, it is very
seldom that the employers complain of this; we might almost say never,
except when it is intended to reduce wages by what might be termed
``working the uniform-rate dodge,''---that is, assuming the rate of the
general run of workmen to be that which should only be paid to the very
best workmen. This is often done when wages are sought to be reduced.
One instance, which will illustrate every other, we select from the
Parliamentary Report on Combinations of Workmen, 1838. Mr. Carolin,
master builder of Dublin, who, through his reduction of wages, was
involved in serious disputes with the workmen, after evading, by a
reference to piece work, the question put by Mr. O'Connell (7394), ``Do
you pay any man more than the minimum stipulated by the body?''
``Admitted (7348) that he paid 27s. a week to his best men and so
downwards.'' By the evidence of Mr. Eaton, also a master builder of
Dublin (7646), the regular wages of the same workmen (carpenters) were
4s. 8d. a day, or 28s. per week, which he paid, ``making no difference
between superior and inferior workmen,'' ``which he thought fair between
master and man,'' but he added, which illustrates what we have stated,
that ``he endeavoured to select the best workmen.'' The one who did not
wish to reduce wages sought and obtained the fair advantage which the
uniform rate afforded; while the other, whose only object was to reduce
wages, complained of this rate, and gave it only to his best hands. And
this is the usual course those who complain of the uniform rate pursue.
The reduction of wages is generally their object, and the uniform rate
is the mere pretence which they use to effect that object in the way
described.

Indeed, so far from placing ``the competent and the incompetent on the
same level,'' this ``uniform rate'' has been bitterly complained of, as
excluding the incompetent altogether. At the late Bradford meeting one
of the speakers gave as a reason against Trades' Unions that as he was
not able to earn the usual rate, and as the Union would not allow any of
its members to work for less, while he was a member he could get no
employment, and so he left it. In ``Chambers' Journal'' for April, 1856,
this complaint is urged with great energy; and Sir J. K. Shuttleworth,
chairman, at the discussion at Bradford on ``Trades' Unions'' in
reference to the above statement, said, ``that nothing was better
attested than that Trades' Unions did coerce their members in this
matter, and that unless they purged themselves from this taint they
would not meet with the sympathy of the public.'' Knowing as we did that
this ``uniform rate of wages,'' which, by the way, is never in any trade
without numerous exceptions, and when the work is done ``by the piece''
does not exist at all; is never in practice a source of difficulty, as
an employer is sure to dismiss the workman who does not suit him---and
we have known excellent workmen so dismissed, as not being used to the
particular kind of work the employer had to execute, though first-rate
hands in their own department; and knowing, also, that this ``uniform
rate'' is rarely complained of by employers, except they wish to reduce
wages, or by workmen unless they are arrant bunglers; we were not a
little surprised to hear of it as a matter of grave complaint. It
becomes therefore necessary to show, in respect to workmen, as we have
previously in respect to employers, how those who thus complain are
affected by this uniform rate, and what kind of workmen they are.

A man may not be worth the regular wages, either because he is slow in
doing his work, or because his work is inferior. In the first case,
where work by the piece is the custom of the trade, there is and can be
no difficulty in the matter, unless, indeed, he be so very slow that no
employer will give him room in his shop. Where, however, work in a trade
is done by the day, there may be a difficulty; still if the man's work
be well executed, he is sure, in the end to get employment where
quantity is not so much regarded in the work as its quality. But if the
man's work be inferior, there is indeed a difficulty, but that
difficulty is not in the regulation of the trade, but in the man
himself. If there were no regulation as to the rate of wages, such a man
would be restricted to those shops only where inferior work was done,
and, if his work were very inferior, even these would be closed against
him; for what master would employ a man to spoil his materials?

This kind of incompetence may proceed from imperfect instruction, in
which case the obvious course for the young man is to engage himself
as an improver; in many cases, however, that we have known, even this
has not been necessary. Where there is a natural aptitude to learn, the
man acquires more skill at every shop he works at, and ultimately
becomes as competent as the best. Where, however, the incompetence from
inferior work is the result of inability to learn, and not for the want
of proper teaching, there is no remedy; indeed, such an one has no
right to designate himself a workman of any trade, if he is not able to
work at it in the style of a workman. But here, again, the difficulty is
not in the regulation, but in the man himself. These, however, are the
men most likely to complain of this regulation. They will lay their
being continually out of employ to any cause rather than to their own
bungling execution of the work entrusted to them to do; to their not
being allowed to work for less than the standard wages; or to anything
but the real cause---their own thorough incompetence.

The reason why all men in a trade are supposed to be competent to earn
and to receive the standard rate of wages is, because, first, the fact
is, that, taken as a whole, they are all competent, the exceptions being
too few to be appreciable; and, second, because, from the predatory
instinct referred to above, which prevails equally among workmen as
among masters, some men, being thoroughly competent, would pretend
incompetence, that they might supplant their fellow workmen by offering
themselves for less than the regular wages. In our own trade, however,
where men from age or obvious inability are unable to earn the regular
wages, there is not, nor ever has been, any objection to their receiving
less.

In reference to the complaint that this uniform rate stints the action
of superior skill or physical strength, such complaint applies, wherever
it justly exists, to the employer as much as to the employed. Workmen
who are remarkably quick are sometimes paid more, but not always, as the
employer mostly likes to reap the benefit of such quickness himself;
which generally results in the quick hand subsiding into the regular
quantity of work. This is often made the ground of great complaint
against workmen that they discourage skill and celerity in their work.
The complaint, however, often should be reversed; the workman might
justly reply that he gets no encouragement. If he were, from rare
physical development, to do double the work of another, he would be paid
no more, or, if he were, it would be by no means in proportion to the
additional work he got through. This is the true reason why workmen are
accused of not wishing to domore than a certain quantity of work in a
day, about which the master builders have been so eloquent and so
denunciatory. The ``Edinburgh and Quarterly Reviews'' have in this
matter taken up their parable against the working man as advocates of
the employers. But why should the workman do more than the regular
quantity of work if he be not paid more? Would either of these two
reviewers write two sheets for the price of one? To expect them to do
more is to desire virtually a reduction of wages, which we have before
shown to be at the bottom of all these complaints of the uniform rate of
wages.

When a trade is unusually brisk, it may be that there are no hands to be
obtained but the inferior workmen, who from this cause have, until then,
remained unemployed; these receiving the usual rate does not alter the
case, as it is with them a question of simple demand and supply, which,
in these circumstances, would determine them to be worth it. Many
complaints admit of this solution, while nothing is said, when things
are exactly the reverse, of superior workmen being obtained---from the
same cause reversed---at the ordinary rate ot wages.

But, after all, we believe both employers and employed are generally
satisfied in respect to this ``uniform rate,'' as it exists in their
respective trades. During a strike, or dispute, each will say hard
things of the other, as men in such circumstances will always do. Apart
from this, if there be men, who will not do the proper quantity of work,
there are exacting employers who, if their men did double the usual
quantity of work, would grudge paying more; but these do not represent
the great body of masters or men in any trade. They are exceptions, of
whom it may justly be said, in reference to their respective
demerits---to use the vernacular---that ``there are six of one, and
half-a-dozen of the other,'' each naturally, and most surely, causing
the result of which they both complain.

But, it is asked, why cannot a man sell his labour for what he likes, as
a shopkeeper tickets his goods under the price of those of his
neighbour? They do not interfere with each other, and why should the
workman? The shopkeepers here alluded to, Ave beg to reply, are not
obliged to be always together, and, therefore, the ill blood, which is
often intense, does not cause interference. But, in wages, the matter
assumes a very different aspect, as will be seen by putting a case.
Suppose, in a shop which employed twelve or any other number of men,
whose wages were 30s. per week, a man was to be engaged who thought
proper, without being asked, to work for 25s. per week; not because he
was deficient in skill or quickness, but simply because he chose to do
so; or, if the work were done by the piece, to offer to do each job at a
corresponding reduction. Would not the rest be very likely to see, as
the result of this, a near prospect of their own wages being lowered to
this standard, and, in consequence, regard the underworker as a
``curry-favour'' and a sneak, who sought unfairly to supplant them, and
treat him accordingly? It may be very true, that this man had a right to
ask what wages he liked for his labour, even to half the above sum; but,
while men are human, Trades' Union or no Trades' Union, he would not be
able to do so without feeling his position to be highly disagreeable to
himself. We should like to know how those who put the above question
would behave if they formed the twelve, or a part of it, above supposed?

But it may be replied, this is exactly the reason why Trades' Unions are
objectionable; they unduly interfere with the natural liberty -which
such a man as this, for instance, undeniably possesses to work for what
he likes. Undoubtedly they do. But the cause of this interference is,
not in the Trades' Union, but in the natural repugnance which fill men
feel to being supplanted, and by such means either turned out of their
employment, or forced to work at a lower wage. It is this repugnance,
which is found everywhere, that prompts the ``undue interference,''
which would be felt whether the men were in combination or not.
Combination might be the effect of this feeling and make it more
effectual, but it is not its cause. Its cause is in the common instincts
of man's nature, which can never be rooted out, nor is it proper that
they should be. To suppose that it is the result of combination is to
mistake the effect for the cause, which indeed, and the assigning of a
wrong cause, are the two great mistakes made by the opponents of Trades'
Unions. The case is different when a reduction of wages proceeds from a
falling market, or from what appears to be a decline in the trade or
manufacture. The natural repugnance to such reduction here gives place
to what is seen to be inevitable. Hence Trades' Unions seldom
interfere---indeed, never---unless through ignorance of the existing
cause. This, however, can hardly happen, because, if the men were ever
so inclined, no strike can take place when there is little or no work to
do. For it will be at once seen, that men who are already unemployed
cannot ``strike'' from their work. When, in the case of the ``Luddites''
of 1812, nearly a whole trade was thrown out of employment by the sudden
introduction of machinery, the disorders which followed were not the
result of any Trades' Union combination, but were riots and violence
impelled by extreme poverty and hunger coming suddenly upon a large
number of people; which, at any time, is likely to recur under the same
circumstances, whether produced by the introduction of machinery, or any
other cause.

\subsection{PIECE WORK.}

It has been alleged that workmen refuse to work by the piece. This is
not correct. There does exist a prejudice against it by some, both of
workmen and employers, but it is not general. The compositors of London,
as a general rule, work by the piece, time-work being the exception;
while those in the country work by the day, the exception being work by
the piece. The shipwrights of the Port of London work by the job or
piece, while those of all other ports work by the day. Indeed, the main
objection in different trades to working by the piece is in the
complaint that, when the men are found to earn good wages at it, the
employer wishes to reduce the price of the work, and that it is so often
made use of, as a means of reducing wages. There are, as we are
informed, in the engineering trades three modes: first, by the piece at
a settled price for the various jobs; second, the giving out a quantity
of work, at a certain price, to one man, who gets others to do it for
less, similar to fhe sweating system among the tailors; and, third, the
selecting of a man who possesses superior physical strength and
quickness as the principal of several workmen, and paying him an
additional rate, by the quarter or otherwise, with the understanding
that he is to exert himself to the utmost to induce the others, who are
only paid the ordinary wages, to keep up to him, by which means is
obtained the work of several men, up to the standard of a workman of
superior working strength and quickness, for the payment, except in the
case of the principal, of the ordinary rate of wages. Without any
comment, this will go far to explain many of the complaints of
``\textit{stinting} the action, superior skill, and working power,''
made by the employers against their men.
    \chapter{---STRIKES.} \label{STRIKES}

When a body of men stand out for a price which their employers refuse to
give, while this dispute is pending, the position of the workmen is that
of a strike. As strikes are the last resort, as they are always
expensive, and as they engender mutual ill feeling, they should never be
entered into without duly calculating the probabilities of success, nor
until all means of amicably settling, the difference have failed. It
often happens that workmen have no alternative, but either to submit to
a reduction tyrannically enforced, without any reasoning on the matter
being allowed, or to cease from labour. Often has a strike thus been
precipitated, and ruin inflicted on employer and employed, which might
have been averted by a little calm reasoning on the matter. It is the
same when a rise in wages is asked by the men. Both parties are apt to
view each other as enemies, and in this jaundiced view, which prevails
equally on both sides---aggravated by the unconciliatory tone which is
sure to result from such a state of feeling---reasoning on the
subject---as the subject, considering the important results to both
paHies which are then pending, should be reasoned upon---is rendered
impossible. The beginning of strife is like the letting out of water
that might be, at the commencement, easily stopped. But if there be one
thing more than another which, in their turn, both parties in these
circumstances often, to all appearance, agree in throwing aside, it is
the conciliatory spirit which might prevent these ruinous disputes. But
while strikes are always to be deprecated, because they are, for the
time, a state of moral warfare, and, like all sates of hostility,
productive of mutual bitterness---and because they are carried on at a
loss to both parties---we are, notwithstanding, clearly of opinion, from
long experience of their results to journeymen both of success and
defeat, that there is no proper alternative, in certain cases, than the
position of a strike.

The following extract from a very talented article in the
\textit{Builder} of August 11, 1859, which very clearly sets forth the
nature of strikes, as clearly exhibits their rationale:---

\begin{quote}
    \hspace{2em} Suppose a question of cotton and sugar value instead of
    labour value---how does the seller know whether he is selling too
    cheap, except by refusing to sell at all below a certain higher
    rate? If there is but little cotton or sugar, as of course he
    suspects or affects to suspect there to be, he will sell what he
    chooses at his own figure; but if not, he must take the buyer's
    price for it. Now, the fixed price about which the transaction halts
    is the strike of the seller against the buyer---of the supplier
    against the demander---and provides the only practicable means of
    arriving at the fair value.
\end{quote}

Against strikes it has been often urged---

\begin{enumerate}
    \item Their great expense;
    \item That they promote the introduction of machinery, and
    consequently leave the workmen in a tenfold worse position than they
    were before.
\end{enumerate}

1. It is admitted that strikes are very expensive. But the expense of
anything must be taken in reference to the gain it is intended to
procure, or the loss it is intended to avert. It should also be viewed
in reference to its result in success or defeat.

The expense of strikes to prevent a reduction of wages, let it be what
it may, is soon equalled by the amount it is intended to take from the
men's wages. A reduction of a penny per hour---we take this sum, because
in our own trade reductions have been attempted which would amount to
it, but any other sum can be taken, as the case may be---in a day often
hours, occurring in a trade of 1,000 men, amounts to £250 a week, or
£13,000 a year.\footnote{At 5 per cent., £13,000 a year represents a
capital of £260,000} This may be called a small trade. In a trade of
10,000 men it amounts to £130,000 a year;\footnote{At 5 per cent.,
130,000 a year represents a capital of £2,600,000} and in a trade of
20,000 men it amounts to £260,000 a year.\footnote{At 5per cent.,
£260,000 a year represents a capital of £5,200,000} Considering the
capital these sums per annum represent, it cannot be surprising that
workmen are willing to incur a very great expense to prevent loss so
enormous.

But suppose the men to be defeated after incurring great expense---the
employers' expenses are sure to be fully equal to those of their men,
besides the possible loss of business. The fact of its being often very
expensive to reduce wages, prevents reductions being attempted which
otherwise would be made without hesitation or scruple. Strikes,
therefore, even in defeat, have a powerful tendency to prevent a future
and further reduction of wages.

Besides, it might as well be urged because in a comparatively short
period millions have been expended in going to law, involving often
total ruin to families, loss of reason, and even suicide, that therefore
no unjust aggressive claim should be opposed, and referred to legal
decision.

With respect to the objection, that strikes promote the introduction of
machinery: at no time is a strike needed, or indeed is there any
stimulus, but the shortening of labour, necessary for its introduction,
whenever it is found to answer. It may be true that ``mules'' and
``double deckers'' were introduced during, or about the time of a
strike; but these, if they saved labour, would have been used if there
had been no strike. The spinning jenny was not invented in consequence
of a strike; neither was the power-loom. No strike existed when the
printing machine was introduced; nor was there any to occasion its
subsequent improvements. Nor has been, or would, the introduction of
machinery in any trade be delayed an hour beyond the time it was found
to answer. ``The ingenuity of our mechanics,'' alluded to in the
``Quarterly Review'' (p. 504) has never needed the impulse of a strike
to call it into action; indeed, the fact has been exactly the contrary.
It has been more apt to employ itself in the making of machines which
were found not to pay, than to require any stimulus, whether of a strike
or otherwise, to call it into action.

Two great mistakes are often committed on both sides in strikes. The
first is, in forgetting that the issue joined is simply to prove which
of the two parties can stand out longest in a bargain---that the dispute
is not contending as in a combat, but simply for the adjustment of a
certain price for labour. It is quite true, that, where a reduction of
wages results from one employer underselling other employers, and who is
seeking to make up his profit out of his men's wages, a sense of wrong
is felt by the men so reduced, similar to that which prompts a man to
repel an injury. And, doubtless, in such a case they think that such an
employer is seeking to extend his trade by taking what does not belong
to him, but to themselves. And, on the other hand, where a rise in wages
is asked, which the employer feels will reduce his profits, and much
narrow his operations, he also will feel the same. Nevertheless, to
bring such a contest to an issue at the least possible cost to both
parties, and the least injury to the trade itself, such feelings ought
to be laid aside, and the dispute conducted according to what it really
is, simply a bargain, in which the employer is the buyer, and the
employed is the seller. Both are standing out against each other's
price. If the seller in this case can do without selling, and the buyer
can get no labour elsewhere, the workman is sure to win. If the
contrary, the employer. It is indispensable, therefore, as the dispute
cannot be carried on without great loss to both, that each party should
be well informed of every circumstance connected with that their true
position, both at the outset and from day to day as the dispute
proceeds. This, however, is not always done, and is the other mistake
often committed. In the builders' strike and ``lock-out,'' if the
masters had at all studied their position as to the number of men likely
to be available in the precipitated strike at Messrs. Trollope's, they
would have found, if their own statements are to be depended on, that
there was no occasion whatever for their lock-out, with its great
expense to themselves, and suffering to their men. It was, as they
expressed it, ``to continue for a month, or until the Messrs. Trollope
and Sons were enabled to resume their operations.''

This firm, by its own statement, was enabled to resume operations in a
fortnight, without having gained over a single man that had been locked
out, either of their own late workmen or those of any other of the
combined employers; and it was the same at the subsequent period, when
they announced that they had obtained their full complement of men. And,
as the lock-out, therefore, did not achieve this, it was indisputably
not needed at all. Their loss of £130,000, as announced in the
\textit{Weekly Dispatch} of January 1, 1860, which, from circumstances
not necessary to allude to here, is unquestionably good authority, has
therefore been incurred without the smallest occasion for it. Many
blunders have no doubt been made by workmen in their strikes, but none
so great as this, because so entirely without real necessity.

Both parties, however, are too often blinded by passion, to see clearly
their way, and being angry, the contest is transformed from what it
really is, merely of a bargain, into that of a combat. Under the
influence of this excitement, the leaders of both parties are too often
more intent upon setting forth the injuries which their respective
audiences imagine they shall sustain, if they abate the least of their
original views; or, what is more mischievous, they feel it a duty to
urge that they cannot without \textit{disgrace} modify, for the purpose
of arrangement, the several points in dispute, forgetting all the while
that it is only a bargain for which they are contending, in which it is
the special business and duty of the seller, if he finds he cannot get
his price, and the market is being, or likely to be, supplied, to take
the next in value, and in no case to stand out, if he sees the market
being supplied without him. Among dealers, he who loses by standing out
too long is always an object of derision. This mistake is too often made
in a strike. The dogged valour which has so often enabled the sons of
toil to conquer in battle is not unfrequently the cause of total defeat
in a strike, when, with a little address, considerable advantage might
have been gained, or loss prevented. Such address has been considered
cowardly and disgraceful, and the affair has ended with the loss of
everything. On the other hand, the employers, under the influence of the
same feelings and mistake, have often been led to refuse all adjustment
or arbitration, and thereby have led, by protracting the struggle, to a
much greater loss than their victory has achieved gain; and in some
cases, as in the master-builders' lock-out, to immense loss, without any
substantial gain at all, unless the gratification of their anger be so
considered.

But, after all, the true position of employer and employed is that of
amity. They are each, notwithstanding these occasional disagreements,
the truest friends of the other, and neither can inflict an injury on
the other without its recoiling on himself. Capital and Labour should go
hand in hand. Experience has amply proved that the Capitalist cannot
injure the Labourer, or the Labourer the Capitalist, without each
inflicting injury, and perhaps ruin, upon themselves.

\section{THE ``EDINBURGH REVIEW.''}

Many of the charges by this reviewer have already been answered. It is
necessary, however, briefly to notice the article. It is written,
apparently, \textit{to order}---confessedly, in profound ignorance of
the subject. This ignorance may be real or feigned---feigned, in order
to induce a belief that Trades' Unions were secret societies, similar to
those which ``honeycomb continental nations,'' abounding in ``spies,''
with ``forced dumbness,'' and ``\textit{surveillance}'' and ``distrust
of neighbours on the right hand and left,'' of which no information
could be with certainty obtained. This, indeed, he largely insinuates,
which is turning his ignorance, if real, to very good account. And it
must be confessed, if this ignorance be feigned, that the acting, if it
be not altogether successful, is very effective. Knowing nothing, or
being obliged, from the course he has adopted, to pretend to know
nothing of the subject, and finding that Adam Smith and McCulloch were
dead against him, and Stuart Mill, who was formerly so, having only
recently modified his opinion and even now enjoining that there should
be no encouragement given to the multiplication of labourers, because in
the opinion of ``farmers,'' ``Boards of Guardians,'' and others of the
employer class, the checking of such multiplication would make them
``too independent,'' and therefore, besides in other respects, not a
presentable authority for his intended article, he with admirable
adroitness tells his readers at the commencement that he is ``not going
to preach political economy.''

Nor was the information respecting Trades' Unions that could be obtained
from ``Blue Books,'' namely, the Parliamentary Enquiry on Combinations of
1838, and the Report of the Hand-loom Weavers' Commission suitable for
his purpose; the last, though quoted by him at the end, not being
applicable to any extent, as the hand-loom weavers were not Trades'
Unionists, which is one reason given in the before-mentioned Enquiry of
1838 (Questions 1418--19) why their wages were so wretchedly low; while
the evidence in the same Enquiry (of 1838), which is divided into two
parts, the first containing chiefly the evidence on the cotton-spinners'
strike of 1837, and the second the evidence on Irish Trades' Unions, is
too conflicting to be used by him, who only intended to see the
employers' side of the question. Indeed, it is so conflicting that the
Committee made no report upon it. Besides, the first of these contained
the opinions of Sir Archibald Alison on Trades' Unions, whom no one will
suspect of having an undue leaning towards them. Apart from intimidation
and violence which we are as much against as he could possibly be he
expresses himself several times favourable to trades' combinations, in
terms similar to the following extract. He is speaking of the
CottonSpinners' Union of Glasgow:---

\begin{quote}
    \hspace{2em} ``Their combination, I think, might have a very
    beneficial effect in their own favour, if it was limited to merely
    legal acts. I think, for example, that it is a very good thing for
    the cotton-spinners, as well as every other class of labour, to
    combine, because it enables numbers, to a certain degree, to
    compensate and to enter with equality into the lists with capital;
    and therefore I think that combinations are essential to support the
    rights of labour in the competition with capital.'' (Question 1956.)
\end{quote}

In this state of things, where it was so essential to shut himself out
from using, if he possessed it, any knowledge of the subject, there was
nothing for it but to follow the course pursued by Lord Chesterfield
when he brought the Bill into the House of Lords for reforming the
Calendar. His lordship tells his son, in one of his Letters, that having
to bring in this Bill; knowing nothing of astronomy, and consequently
unable to explain the reasons for the change; he did what he believed
his audience would be better pleased with, namely, instead of urging the
reason for the alteration from the science---which, he further said, he
did not believe his audience would have understood if he had been able
to give such explanation---he gave a history of the changes made in the
Calendar down to the then present time, taking care to flatter the
prejudices of his hearers by telling them what he thought they would
like to hear, by which they were so delighted that they gave him credit
for his profound knowledge of the science; while another noble lord, who
well understood astronomy, and who handled the subject with the learning
of a philosopher, was heard with impatience. With this example of the
advantage of being ignorant of the subject upon which he was about to
write, which he betters, by turning his real or assumed ignorance to the
account described:---he proceeds to give a history of alleged misdeeds
of workmen when combined---of many of their strikes, in which they are
described to be everything that is unreasonable, unjust, tyrannical, and
absurd, in the hope, as already in part quoted, that his readers will
believe of Trades' Unions that---

\begin{quote}
    \hspace{2em} Their aim and object is, in every case which we have
    been enabled to investigate, to \textit{stint} the action of
    superior physical strength, moral industry, or intelligent skill; to
    depress the best workman in order to protect the inferior workman
    from competition; to create barriers which no Society-man can
    surmount, and which few non- Society men dare to assail; and, in
    short, to apply all the fallacies of the protective system to
    labour. Such a system injures, first, the individual, whom it robs
    of a free market for his labour; secondly, the class of manufactures
    to which he belongs, by increasing the cost and diminishing the
    efficiency of the workmen; and, lastly, the nation at large, by
    curtailing the productive power, and consequently the wealth of the
    community.---P. 529.
\end{quote}

The only result of his ``investigation,'' as exhibited by himself, has
been the total exclusion of every fact and circumstance that would
justify, excuse, or extenuate the proceedings of the men, which, indeed,
appears to have been the only object he had in view in setting himself
to investigate the matter.

We have already shown the fallacies of the assertions, upon which this
confused argument rests, itself consisting principally of mere
assertion. The only part we have not alluded to that might be construed
into an application of the ``fallacies of the protective system to
labour,'' is the insisting, as Trades' Unions do, that all their members
should serve an apprenticeship, generally of seven years, to their
trade, and of sometimes endeavouring to limit the number of apprentices.
This last, however, is not now the general practice, but the exception.
It is always a needy or grasping employer who takes an inordinate number
of apprentices, and he never cares a straw for ``the nation at large,''
or ``the wealth of the community;'' he only aims, in thus using ``its
productive power,'' to enrich himself, no matter at whose expense.

The Act of Elizabeth, requiring an apprenticeship of seven years before
exercising a trade, has been repealed nearly fifty years, after having
been in force for two centuries and a half The ``Quarterly Review,'' in
its late article against Trades' Unions, states that this Act
``unquestionably exercised an important influence on English industry.
It stigmatised and punished the idle and the vagabond, directed the mass
of the people to manual occupation as affording the best means of
independent subsistence; and being acted on steadily from generation to
generation, it gradually educated a nation of skilled artizans.''---P.
480. And although the Act has been repealed so long, the practice has
not fallen into disuse. In some trades there is, probably, less to be
learned than in others; but there are very few trades in which long
experience and practice are not necessary, before the workman is able to
compete with those who are considered valuable to the employer. This
ability must be acquired at some period, and what more proper than that
in which the individual passes from youth to manhood---a time in which,
of all others, a facility of hand, and dexterity in performing difficult
processes, is the easiest acquired. This, then, being the most proper
period, and at the same time this also embracing the period when the
passions develope themselves, and are the most ungovernable, does it not
follow irresistibly that, if ever in a man's life he need the
experienced direction and restraint of his elders, it is this period? An
apprenticeship of seven years, from the age of fourteen to twenty-one,
whether considered necessary to learning a trade, as the best period for
acquiring it, or as conducing to the future welfare of the apprentice
himself, points itself out therefore as absolutely necessary. And this
is the reason why the repeal of the Act has made so little difference in
the practice, and why Trades' Unions, generally speaking, insist that
all their members should have served ``their time as apprentices.'' The
members are most of them fathers, and they know by experience the
beneficial effects of the practice, and the mischief that attends the
letting boys ``do as they like'' at this critical period. But we shall
be told this is no reason why the ``productive power, and consequently
the wealth of the community,'' should be ``curtailed.'' True. And the
same reason exists why the ``master villany of the earth, which includes
every other villany,'' and deadly sin---slavery, should perpetually
continue; for how else can ``the wealth of the community'' be so well
kept up by supplies of cheap cotton and sugar?

By the ``nation at large'' we suppose is meant what is termed the
``public.'' The public cares very little how things are produced, so it
gets them cheap, and the cheaper they become the more the public likes
it. The public, however, always respects those who so take care of
themselves as to prevent their suffering wrong. All the reasoning in the
world, nor all the admonitions of virtue, benevolence, or religion,
would prevent the public from preferring slave-grown sugar or
cotton---all the horrors of the slave trade notwithstanding---if it were
cheaper and as good as either grown by free labour. But, in justice to
the public it must be admitted, that it is also true, that, in the event
of men refusing to be slaves---supposing such a thing were to
happen---it would receive with a shout of mingled execration and
derision the wailing of those who complained that commodities were dear
because men were no longer enslaved to make them cheap. The reviewer,
therefore, has no occasion to make himself uneasy about the ``nation at
large;'' it never scruples to take care of itself, or fails to respect
those who do the same.

In giving his history of the alleged misdeeds of Trades' Unions, he has
not bettered his example. Nothing was ever overdone by Lord
Chesterfield. Indeed, on second thoughts we begin to doubt whether he
ever read the works of that nobleman. He proves too much by his
exaggeration; that is, if a jumble of alleged facts, disjointed from all
proof that they of necessity belong to Trades' Unions, can be said to
prove anything. And it is superfluous to say that he who proves too much
proves nothing. He wishes it to be believed that these misdeeds were
done by Trades' Unionists in and by their Union. If even they were all
true, to make them prove anything against Trade Societies it is
indispensable that he should also show of each and all that Trades'
Unions could not exist without these or similar enormities. This he does
not attempt to do; indeed he has entirely shut himself out from doing
so, by his real or assumed ignorance of their action---by his own
representation that they are secret, and that therefore nothing can be
known about them. He seems to have forgotten that his real or pretended
ignorance of their nature and action, so convenient for dark
insinuation, makes him incompetent even to argue the subject, much less
to prove anything, and consequently worthless as an authority. We have
already shown the utter fallacy of his insinuation, that Trade Societies
are secret, irresponsible in their government, or tyrannical; but that,
on the contrary, they are instituted for the purpose of correcting the
disadvantage inherent to the position of the working man when he
bargains \textit{singly} for the sale of his labour, namely, of being
compelled by his immediate necessities to take a price for it below its
fair exchangeable value. To assert that they, in any given instances,
have been turned to a bad purpose, is only to say of them what may with
truth be said of everything human.

The article consists of thirty-eight pages, out of which twenty-eight
are devoted to this history. Some fourteen or sixteen cases are
mentioned, which we suppose are by him deemed sufficient proof; and if
the careful suppression of everything alleged by the men as the reason
of their proceedings, with the aggravation of all that was alleged
against them, and the as careful setting forth as patterns of justness
and liberality the proceedings of the employers, could achieve a
demonstration of truth, the reviewer would be triumphant. It is rare to
find so much alacrity, so much \textit{con amore} readiness to accuse.
If the reviewer had been formerly an Unionist---a compositor, for
instance---and, like Hugh Miller, had been affronted by his former
associates, he could not be more bitter against Trades' Unions. His zeal
against them is as great and as indiscreet as that of an apostate. The
most absurd and exorbitant demands, upon the silliest pretences, the
most unjust and tyrannical proceedings towards each other, without sense
or reason, he charges, in the various instances he has given, against
the members of Trades' Unions. He refutes himself from the sheer
impossibility of such allegations to be true; or, which is the same
thing, to contain the whole truth. Our argument does not rest upon the
refutation of particular cases; it is, therefore, not necessary to
follow him through his twenty-eight pages of one-sided statement. It is
only necessary to take the first-mentioned, which is the strike of the
tin-plate workers, at Mr. Perry's, Wolverhampton, by which may be judged
how far he is to be relied upon for the rest. He revels in bitter
accusation with the vehemence of female malignity, incapable, like a
true virago, of seeing wTong in anything or anybody but the object of
her vituperation, recounting each accusation without pause or break in
the energy of her volubility, until her vehemence has fairly spent
itself. Such is the manner in which the subject is treated, and, as
might be expected, the conclusion is destitute of logical sequence.

After a page or two of uncertain maundering on the Hand-loom Weavers'
Commission, which he quotes, and with which he does, and does not, agree
in respect to alterations in the law of combination suggested therein,
which he very truly suggests ``do not amount to much,'' the conclusion
he comes to is, that there should be a parliamentary inquiry into the
state of industrial society in England, before a select committee. What,
an ``inquiry'' only after so much violent accusation about the
``mysterious tyranny,'' which has ``honeycombed'' English
society---eaten, as it were, into its very vitals---like that of
Continental nations, with its ``spies,'' ``enforced dumbness,'' "
\textit{suveillance}" and ``distrust on the right hand and on the
left?'' Why, he abjures himself. It is a confession that he has written
to the prejudices of his readers, and that he is frightened lest they
should seriously believe what he himself does not know to be true.
Though this is quite apparent from his mode of treating the subject, it
is singular that he should virtually confess to the fact. Yet with
strange perversity, like a virago, who, no matter how she has destroyed
her own credibility, ever returns to her first word, he reiterates as
his last word, this time by an assertion, in answer to a roundabout
question put by himself, that Trades' Unions are ``a secret organisation
of trades,'' which ``have undermined the groundwork of society in
England!''

We now proceed to show how far his version of the first of the cases he
says he ``investigated'' is to be depended on, by giving an account of
the same case, chiefly from the same authorities, which he himself
consulted in his investigation. The minutes quoted of the Committee of
the National Association, to which he had not access, only showing the
success with which the defendants were imposed upon, of the imposition
itself there was ample evidence in the report of the trial, which the
reviewer states he consulted.

Four pages are devoted to an account of the strike and subsequent trial
for conspiracy at Wolverhampton. The strike was at Mr. E. Perry's,
tin-plate worker, japanner, \&c., of that place, who was afterwards the
prosecutor.

Mr. Perry is said to ``have believed'' that he was paying ``ten per
cent, higher wages than the average of his trade in Wolverhampton; and
there seemed to be no conceivable reason why anybody should be troubled
about the affairs of the manufactory.'' According to the men, however,
Mr. Perry had been for years attacking their wages; and, at the time of
the strike, so far from his paying 10 per cent, more, he was paying 15
per cent, less than the average wages of the other manufacturers of the
place.

Mr. Perry admitted, in his evidence at the trial, that there had been a
strike in his factory in 1842; and, after some hesitation, that he had
agreed to a ``particular'' book of prices for his own shop, and that
``he had been distinguished in Wolverhampton by having a great deal of
opposition brought against his manufactory, by other persons not
connected with it.'' This does not much differ from the men's account of
the ``opposition'' to which he had been subjected through paying lower
wages than the other manufacturers of the town. With certain employers,
it is never their own men that are dissatisfied; it is always other
people who set them on. To rectify this, the men joined the ``National
Association for the Protection of Industry,'' and requested its
Committee to wait upon Mr. Perry in their behalf. Accordingly, Messrs.
Peel and Green, and afterwards Messrs. Green and Winters, in company
with ``three or four of his own men,'' waited upon him with a list of
prices, to which they wished him to agree. One of this deputation
afterwards became his foreman, and was one of the witnesses against the
defendants. The other employers, six firms in all, were also waited
upon, with the same list of prices, two of whom---Mr. W. E. Walton and
Messrs. Shoolbred, who were large manufacturers, and who employed more
men than all the others put together---returned a favourable answer;
one, Mr. Thurston, a conditional acceptance; the other three, Mr. E.
Perry, Mr. R. Perry, and Mr. Fearncomb, entered into a negotiation with
the deputation, or professed to do so. In reference thereto, Mr. E.
Perry, in a letter to the \textit{Daily News}, of August 20, 1851,
writes:---``I gave the parties to understand that I required time to
examine their proposal, and more especially to look into the book of
prices.''

In furtherance of this expressed intention, and to keep up this
appearance, Mr. E. Perry on several occasions met the defendants, and
actually agreed to the price of certain articles; at which time there
appeared to be the most perfect accord as to the negotiation. Mr. Peel,
as admitted by Mr. E. Perry, said that ``he was not for strikes, but for
peace,'' to which Mr. E. Perry replied, that ``if a satisfactory
arrangement took place, it would be through his (Mr. Peel's) good
management.'' Mr E. Perry also drew up a preliminary agreement
declaratory of such arrangement, which they that is, Mr. E. Perry, Mr.
Peel, and the rest of the defendants who were present, except
Rowlands---all signed.

How far the members of the Committee of the National Association were
justified in believing themselves, besides being appointed by the
journeymen tin-plate workmen of Wolverhampton, to represent them, also
fully recognised in that capacity by Mr. E. Perry, will appear by the
following extracts from the minute-book, written at the time, of the
National Association:

\begin{quote}
    \hspace{2em} ``April 9, 1850. At an interview that Messrs. Peel and
    Green had with Mr. Perry, it was agreed, in consequence of the
    meeting of the ironmasters then about to take place, to postpone all
    further proceedings for a fortnight; but still it was deemed
    necessary that the other employers should be personally visited, 
    Mr. E. Perry, in the meantime, promising to use his influence to obtain
    a conference with the other employers on the subject.''

    \hspace{2em} ``May 3, 1850. Mr. Green reported that Mr. E. Perry had
    consented to a conference of the employers and the men, one from
    each shop, together with Messrs. Peel and Green, as their advisers
    and advocates, to arrange the book of prices.''

    \hspace{2em} ``May 30, 1850. Mr. Green reported that Mr. Peel and
    himself had had two meetings with Mr. E. Perry about the
    preliminaries, and that the conference of the men and the employers
    would take place on Monday, June 3, 1850.''

    \hspace{2em} ``June 5, 1850. Mr. Green reported that the proposed
    conference between the employers and the men, commenced on Monday;
    that there were four employers and six men (one from each shop)
    present, in addition to himself and Mr. Peel, to arrange the book of
    prices.''
\end{quote}


It was at this meeting that Mr. E. Perry drew up the preliminary
agreement declaratory of such arrangement referred to above.

We give these extracts, not for the purpose of comment, but simply to
show that, whatever might be Mr. E. Perry's version of the affair, the
two members of the Committee of the National Association had good reason
to believe that he was perfectly willing to arrange respecting the
prices of the men. We may here mention, that, in the end, the two
employers, Messrs. Shoolbred and Co., and Messrs. W. E. Walton and Co.,
whom Mr. E. Perry, in his letter of August 20, calls the two ``youngest
firms,'' who, as before observed, employed more men than all the other
firms of the place put together, agreed to the terms of the men.

\end{document}
